\subsection{与李代数的正则元的关系}

命题 6. 设 V 是 g 的一个开子群，并设 exp : V → G 是定义在 V 上的一个指数映射。

\begin{enumerate}
\def\labelenumi{(\roman{enumi})}
\item
  存在一个邻域 \(\mathrm{W}\)（0 在 \(\mathrm{V}\) 中的邻域），使得 \(\mathfrak{g}^1(\exp x) = \mathfrak{g}^0(x)\) 对所有 \(x \in \mathrm{W}\) 成立。
\item
  若 \(k = \mathbf{R}\) 或 \(\mathbf{C}\)，则 \(\mathfrak{g}^1(\exp x) \supset \mathfrak{g}^0(x)\) 对所有 \(x \in \mathfrak{g}\) 成立。
\end{enumerate}

由第三章 §4 第 4 节命题 8 的推论 3，存在一个邻域 \(V'\)（0 在 \(V\) 中的邻域），使得对所有 \(x \in V'\)，\(\exp(\mathrm{ad}(x)) = \sum_{n=0}^{\infty} \frac{1}{n!} \mathrm{ad}(x)^n\) 有定义且 \(\mathrm{Ad}(\exp(x)) = \exp(\mathrm{ad}(x))\)。若 \(P \in k[X]\) 且 \(\alpha \in \mathrm{End}(\mathfrak{g})\)，容易验证 \(\mathfrak{g}^\lambda(\alpha) \subset \mathfrak{g}^{P(\lambda)}(P(\alpha))\) 对所有 \(\lambda \in k\) 成立。因而，

\[
\mathfrak {g} ^ {0} (\operatorname{ad} (x)) \subset \mathfrak {g} ^ {1} (\exp (\operatorname{ad} (x))) = \mathfrak {g} ^ {1} (\operatorname{Ad} (\exp x)) = \mathfrak {g} ^ {1} (\exp x)
\]

对所有 \(x \in V'\) 成立。若 k = R 或 C，则 V = g，且可取 \(V' = V\)，这就证明了 (ii)。我们证明 (i)。设 U 是 \(\text{End}(\mathfrak{g})\) 中 0 的一个邻域，使得 \(\text{Log}(1 + \alpha) = \sum_{n > 0} (-1)^{n+1} \frac{1}{n} \alpha^n\) 对所有 \(\alpha \in U\) 有定义。于是在 0 的一个邻域上 \(\log \circ \exp = 1\)，且 \(\mathfrak{g}^1(1 + \alpha) \subset \mathfrak{g}^0 (\text{Log}(1 + \alpha))\) 对所有 \(\alpha \in U\) 成立。设 W 是 g 中 0 的这样一个邻域：它由那些 \(x \in V'\)（满足 \(\exp x \in 1 + U\)）组成，并且

\[
\operatorname{Log} (\exp (\operatorname{ad} (x))) = \operatorname{ad} (x).
\]

于是，对所有 \(x \in W\)，

\[
\begin{array}{c} \mathfrak {g} ^ {1} (\exp x) = \mathfrak {g} ^ {1} (\mathrm{Ad} (\exp x)) = \mathfrak {g} ^ {1} (\exp (\mathrm{ad} (x))) \\ \qquad \subset \mathfrak {g} ^ {0} (\operatorname{Log} (\exp (\operatorname{ad} (x)))) = \mathfrak {g} ^ {0} (\operatorname{ad} (x)) = \mathfrak {g} ^ {0} (x). \end{array}
\]

这表明 \(\mathfrak{g}^1 (\exp x) = \mathfrak{g}^0 (x)\) 对所有 \(x\in \mathrm{W}\) 成立。

引理 4. 设 U 是 g 中 0 的一个邻域，exp 是 U 到 G 的一个指数映射，它在 U 的每一点处平展，且使 \(\mathfrak{g}^{1}(\exp x) = \mathfrak{g}^{0}(x)\) 对所有 \(x \in U\) 成立。

\begin{enumerate}
\def\labelenumi{(\roman{enumi})}
\item
  函数 \(r_{\mathrm{Ad}}^0\) 是常数且等于 \(\mathfrak{g}\) 的秩（在 \(\exp(\mathrm{U})\) 上）。
\item
  若 \(x \in \mathrm{U}\)，则 exp \(x\) 正则当且仅当 \(x\) 是 \(\mathfrak{g}\) 的正则元。
\item
  元素 \(a \in \exp(\mathrm{U})\) 正则当且仅当 \(\mathfrak{g}^1(a)\) 是 \(\mathfrak{g}\) 的一个嘉当子代数。
\end{enumerate}

设 \(l = \operatorname{rk}(\mathfrak{g})\)。若 \(x \in \mathrm{U}\) 是 \(\mathfrak{g}\) 的一个正则元，

\[
r _ {\mathrm{Ad}} (\exp x) = \dim \mathfrak {g} ^ {1} (\exp x) = \dim \mathfrak {g} ^ {0} (x) = l.
\]

由于属于 U 的 g 的正则元构成 x 的一个邻域，且 exp 在 x 处平展，这表明 exp x 正则且 \(r_{\mathrm{Ad}}^{0}(\exp x) = l\)。由于属于 U 的 g 的正则元在 U 中稠密，我们有 \(r_{\mathrm{Ad}}^{0}(a) = l\) 对所有 \(a \in \exp(\mathrm{U})\) 成立。设 \(a \in \exp(\mathrm{U})\) 是 G 的一个正则元，并设 \(x \in U\) 使得 \(a = \exp x\)。由于 \(\mathfrak{g}^{0}(x) = \mathfrak{g}^{1}(a)\)，有 \(\dim \mathfrak{g}^{0}(x) = r_{\mathrm{Ad}}^{0}(a) = l\)。于是，x 是 g 的正则元，且 \(\mathfrak{g}^{1}(a)\) 是 g 的一个嘉当子代数。最后，若 \(a \in \exp(\mathrm{U})\) 且 \(\mathfrak{g}^{1}(a)\) 是 g 的一个嘉当子代数，

\[
r _ {\mathrm{Ad}} (a) = \dim \mathfrak {g} ^ {1} (a) = l = r _ {\mathrm{Ad}} ^ {0} (a),
\]

于是 \(a\) 正则。

命题 7. 设 V 是 e 在 G 中的一个邻域。g 的每个嘉当子代数都形如 \(\mathfrak{g}^{1}(a)\)，其中 a 是属于 V 的 G 的正则元。

由命题 6，存在 g 中 0 的一个开邻域 U 和一个指数映射 \(\exp: U \to G\)，满足引理 4 的条件。若 h 是 g 的一个嘉当子代数，则存在正则元 \(x \in h\) 使得 \(\mathfrak{h} = \mathfrak{g}^{0}(x)\)（§3，定理 2）。另一方面，存在元素 \(t \in k^{*}\) 使得 \(tx \in U\) 且 \(\exp(tx) \in V\)。于是 \(\mathfrak{h} = \mathfrak{g}^{0}(x) = \mathfrak{g}^{0}(tx) = \mathfrak{g}^{1}(\exp(tx))\)，且由引理 4 (ii)，\(\exp(tx)\) 是 G 的一个正则元。

命题 8. 设 \(l\) 是 \(\mathfrak{g}\) 的秩。存在一个开子群 \(G^*\)（\(G\) 的开子群），使得：

\begin{enumerate}
\def\labelenumi{(\roman{enumi})}
\item
  函数 \(r_{\mathrm{Ad}}^0\) 在 \(\mathbf{G}^*\) 上是常数，其值为 \(l\)；
\item
  元素 \(a \in \mathbf{G}^*\) 正则当且仅当 \(\mathfrak{g}^1(a)\) 是 \(\mathfrak{g}\) 的一个嘉当子代数；
\item
  若 \(a \in \mathbf{G}^*\)，则 \(\mathfrak{g}^1(a)\) 的每个嘉当子代数都是 \(\mathfrak{g}\) 的嘉当子代数。
\item
  由命题 6，存在 g 中 0 的一个开邻域 U 和一个从 U 到 G 的指数映射 exp，满足引理 4 的条件。以下用 \(G^{*}\) 表示 G 的单位元连通分支（若 k = R 或 C），或 G 的一个包含于 \(\exp(\mathrm{U})\) 的开子群（若 k 是超度量的）。由于 \(r_{Ad}^{0}\) 局部常值且它在 \(\exp(\mathrm{U})\) 的任一点的值都是 l（引理 4 (i)），推出 \(r_{Ad}^{0}\) 在 \(G^{*}\) 上是常数且等于 l。
\item
  设 \(\mathrm{R}^*\)（相应地，\(\mathrm{S}^*\)）是 \(\mathrm{G}^*\) 的正则元所成的集（相应地，元素 \(a \in \mathrm{G}^*\)（使 \(\mathfrak{g}^1(a)\) 是 \(\mathfrak{g}\) 的嘉当子代数者）所成的集）。则 \(\mathrm{S}^* \subset \mathrm{R}^*\)。事实上，若 \(a \in \mathrm{S}^*\)，则 \(r_{\mathrm{Ad}}(a) = l = r_{\mathrm{Ad}}^0(a)\)。我们证明 \(\mathrm{R}^* \subset \mathrm{S}^*\)。若 \(k\) 是超度量的，这由包含关系 \(G^{*} \subset \exp(U)\) 与引理 4 (iii) 得出。设 k = C。由命题 5 的推论，若 \(a \in R^{*}\)，则对属于 a 的一个邻域的每个 \(a'\)，\(\mathfrak{g}^{1}(a')\) 经 g 的一个自同构与 \(\mathfrak{g}^{1}(a)\) 共轭。这证明 \(S^{*}\) 与 \(R^{*} - S^{*}\) 都是 \(G^{*}\) 的开子集。我们已看到 \(S^{*}\) 包含 e 的一个邻域中的所有正则元（引理 4 (iii)）；因而 \(S^{*}\) 非空。由于 \(G^{*}\) 连通，\(R^{*}\) 也连通（命题 1），从而 \(S^{*} = R^{*}\)。
\end{enumerate}

剩下研究情形 \(k = \mathbf{R}\)。首先设 \(\mathbf{G}^*\) 是 \(\mathbf{GL}(\mathrm{E})\) 的一个积分子群，其中 \(\mathrm{E}\) 表示一个有限维实向量空间。设 \(\mathrm{G}_c^*\) 是 \(\mathbf{GL}(\mathrm{E} \otimes_{\mathbf{R}} \mathbf{C})\) 的积分子群，其李代数为 \(\mathfrak{g}_c = \mathfrak{g} \otimes \mathbf{C}\)。存在 \(\mathrm{G}_c^*\) 上的一个解析函数，其零点集是 \(\mathrm{G}_c^*\) 的正则元素所成开集的余集。由《微分与解析流形·结果》3.2.5，这个函数不能在 \(\mathrm{G}^*\) 的每一点为零。因而 \(\mathrm{G}^*\) 含有 \(\mathrm{G}_c^*\) 的一个正则元。设 \(\mathrm{Ad}_c\) 是 \(\mathrm{G}_c^*\) 的伴随表示。对任何 \(a \in \mathrm{G}^*\)，\(\mathfrak{g}_c^1(a) = \mathfrak{g}^1(a) \otimes \mathbf{C}\)，故 \(r_{\mathrm{Ad}_c}(a) = r_{\mathrm{Ad}}(a)\)。若 \(a \in \mathrm{G}^*\) 是 \(\mathrm{G}_c^*\) 的一个正则元，则它是 \(\mathrm{G}^*\) 的正则元，且 \(r_{\mathrm{Ad}_c}^0(a) = r_{\mathrm{Ad}}^0(a)\)。由于函数 \(r_{\mathrm{Ad}_c}^0\) 与 \(r_{\mathrm{Ad}}^0\) 分别在 \(\mathrm{G}_c^*\) 与 \(\mathrm{G}^*\) 上是常数，由此推出 \(\mathrm{G}^*\) 的正则元就是 \(\mathrm{G}_c^*\) 的那些属于 \(\mathrm{G}^*\) 的正则元。由上述，若 \(a\) 是 \(\mathrm{G}^*\) 的一个正则元，则 \(\mathfrak{g}_c^1(a) = \mathfrak{g}^1(a) \otimes \mathbf{C}\) 是 \(\mathfrak{g}_c\) 的一个嘉当子代数；这蕴含 \(\mathfrak{g}^1(a)\) 是 \(\mathfrak{g}\) 的一个嘉当子代数（§2，命题 3）。

现在设 G 是单连通的。存在一个有限维实向量空间 E 和一个从 G 到 \(G'\)（\(\mathbf{GL}(\mathrm{E})\) 的一个积分子群）的平展态射 f（第三章，§6，第 1 节，定理 1 的推论）。由命题 2，若 \(a \in G\) 正则，则 \(f(a)\) 正则。由前所述，\(\mathfrak{g}'^{1}(f(a))\) 是李代数 \(g'\)（\(G'\) 的李代数）的一个嘉当子代数。由于 \(\mathfrak{g}'^{1}(f(a)) = (\mathrm{T}f)\mathfrak{g}^{1}(a)\) 且 Tf 是 g 到 \(g'\) 的同构，这就证明 \(\mathfrak{g}^{1}(a)\) 是 g 的一个嘉当子代数。

最后转向一般情形 \((k = \mathbf{R})\)。设 \(\tilde{G}\) 是 \(G^{*}\) 的一个泛覆盖，\(\tilde{\mathfrak{g}} = \mathrm{L}(\tilde{\mathrm{G}})\)，q 是 \(\tilde{G}\) 到 \(G^{*}\) 的典则映射。由于 q 的核包含于 \(\tilde{G}\) 的中心，若 \(a \in G^{*}\) 正则且 \(a' \in q^{-1}(a)\)，则 \(a'\) 正则（命题 2）。由前所述，\(\tilde{\mathfrak{g}}^{1}(a')\) 是 \(\tilde{g}\) 的一个嘉当子代数。由于 \(\mathfrak{g}^{1}(a) = (\mathrm{T}q)\tilde{\mathfrak{g}}^{1}(a')\) 且 Tq 是 \(\tilde{g}\) 到 g 的同构，这就证明 \(\mathfrak{g}^{1}(a)\) 是 g 的一个嘉当子代数。

\begin{enumerate}
\def\labelenumi{(\roman{enumi})}
\setcounter{enumi}{2}
\tightlist
\item
  由命题 5，存在一个邻域 \(\mathrm{V}\)（\(a\) 的邻域），使得对所有 \(a' \in \mathrm{V}\)，\(\mathfrak{g}^1(a')\) 与 \(\mathfrak{g}^1(a)\) 的一个子代数经 \(\mathfrak{g}\) 的一个自同构共轭。由于 \(a\) 的每个邻域都含有 \(\mathbf{G}^*\) 的一个正则元，由 (ii) 推出 \(\mathfrak{g}^1(a)\) 含有 \(\mathfrak{g}\) 的一个嘉当子代数。于是，由 §3 的命题 3，\(\mathfrak{g}^1(a)\) 的每个嘉当子代数都是 \(\mathfrak{g}\) 的嘉当子代数。
\end{enumerate}

注记. 若 \(k = \mathbf{C}\)，则子代数 \(\mathfrak{g}^1(a)\)（\(a\) 正则且属于一个连通分支 \(M\)（\(G\) 的连通分支））在 \(\operatorname{Int}(\mathfrak{g})\) 下相互共轭。事实上，设 \(R\) 是 \(G\) 的正则元所成的集。对所有 \(a \in R \cap M\)，设 \(M_a\) 是那些 \(b \in R \cap M\)（使 \(\mathfrak{g}^1(a)\) 与 \(\mathfrak{g}^1(a)\) 在 \(\operatorname{Int}(\mathfrak{g})\) 下共轭者）所成的集。我们有 \(\operatorname{Int}(\mathfrak{g}) = \operatorname{Ad}(G^0)\)，其中 \(G^0\) 是 \(G\) 的单位元连通分支。由命题 5 的推论，\(M_{a}\) 在 R 中是开的。由此推出 \(M_{a}\) 在 R 中既开又闭。由于 \(k = C\)，\(R \cap M\) 连通（引理 1），故 \(M_{a} = R \cap M\)。

